\documentclass[10pt,a4paper]{amsart}
\usepackage[margin=19mm]{geometry}
\usepackage{amsmath,amssymb,mathtools}
\usepackage[scr=rsfs,cal=euler]{mathalfa}
\usepackage{xfrac}
\usepackage[unicode,hidelinks,bookmarks=false]{hyperref}
\newcommand{\ie}{i.e.\ }
\newcommand{\cf}{cf.\ }
\newcommand{\resp}{respectively\ }
\newcommand{\Subsection}{\subsection}
\providecommand{\nobreakdash}{\nobreak\hbox{-}}
\setlength{\emergencystretch}{2em}
\multlinegap0pt
\usepackage{bm}
\usepackage{bbm}
\usepackage{dsfont}
\usepackage{xspace}
\usepackage{cancel}
\usepackage{mathtools}
\usepackage{amsthm}
\makeatletter
\@ifundefined{theorem}{\newtheorem{theorem}{Theorem}}{}
\makeatother
\title[Quasi-rectifiable Lie algebras for partial differential equa]{Quasi-rectifiable Lie algebras for partial differential equations}
\author{A.M. Grundland, J. de Lucas}
\date{Source: arXiv:2312.05238v2 (CC BY 4.0)}
\begin{document}
\maketitle

\noindent\emph{Source and licence.} Quasi-rectifiable Lie algebras for partial differential equations. Version arXiv:2312.05238v2 (CC BY 4.0), distributed under the Creative Commons Attribution 4.0 International licence, as recorded in the arXiv licence metadata for this exact version and shown on the abstract page https://arxiv.org/abs/2312.05238v2. Full rights notice: licenses/arxiv-2312.05238-CC-BY-4.0.txt.\par\smallskip

\noindent\textit{Editorial scope.} Counted unit: the Theorem whose source label is \texttt{Th:HamRec} in the article (source0.tex lines 560--561, source label \texttt{Th:HamRec}), delivered together with the article's own complete original proof of it (source0.tex lines 562--600, one \texttt{proof} environment of 39 lines). No mathematics is written, repaired or re-derived: statement and proof are the article's own text line for line. Required context retained uncounted: none. Every definition, display and label of those lines is retained unchanged. Omitted from this package and not redistributed: 1--559, 601--1736. Nothing inside a retained excerpt is omitted or altered: every retained body is delivered exactly as the article's lines are written, so no omission receipt of any kind accompanies this package. Citations: the retained excerpts carry no citation command at all, so no bibliography entry of the article is transcribed and no reference is redistributed. Reference adaptation: none was needed and none was made; every reference inside the delivered unit resolves inside it, so no unresolved reference or citation is produced. Printed slips kept literal: no printed word, symbol, hypothesis or number of the counted statement or of its proof was altered. Layout and class: the article's own class is \texttt{article} with packages including inputenc, physics, parskip, lscape, hyperref, tcolorbox, caption, amsmath, amsthm, amssymb, relsize, mathrsfs, braket, mathtools; the wrapper is the lane's pinned \texttt{amsart} editorial wrapper at 10pt on A4, which restates the theorem-like environments the retained text needs and adds the article's own macro definitions that the retained text needs (none), with extra wrapper packages (bm, bbm, dsfont, xspace, cancel, mathtools). The delivered numbering is the unit's own. Assets: no retained span contains a figure environment, a graphics-inclusion command, a PDF-page inclusion, a table or a TikZ picture; no image file is copied into this package, the package declares no asset dependency and no image file is redistributed with it. Licence: version arXiv:2312.05238v2 of this article is distributed under the Creative Commons Attribution 4.0 International licence, as recorded in the arXiv licence metadata for this exact version and shown on the abstract page https://arxiv.org/abs/2312.05238v2. A full rights notice is delivered at licenses/arxiv-2312.05238-CC-BY-4.0.txt. Characters: every retained excerpt is pure ASCII, so the delivered engine, XeLaTeX with its default Latin Modern text font, reports no missing character.
\begin{theorem}\label{Th:HamRec} A family of vector fields $X_1,\ldots,X_r$ on $N$ is a quasi-rectifiable family of Hamiltonian vector fields relative to a symplectic form $\omega$ on $N$ if and only if there exists a family of Hamiltonian functions $h_1,\ldots,h_r$ on $N$ for the vector fields $X_1,\ldots,X_r$, respectively, such that the Poisson bracket of $h_i,h_j$ is of the form $\{h_i,h_j\}=h_{ij}(h_i,h_j)$ for some functions $h_{ij}\in C^\infty(\mathbb{R}^2)$ with $1\leq i<j\leq r$.
\end{theorem}

\begin{proof} Since $X_1,\ldots,X_r$ form a quasi-rectifiable family of vector fields, one has  $[X_i,X_j]=f_{ij}^iX_i+f_{ij}^jX_j$ for $1\leq i<j\leq r$ for certain functions $f_{ij}^i,f_{ij}^j\in C^\infty(N)$. If, in addition, $X_1,\ldots,X_r$ are Hamiltonian vector fields relative to $\omega$, then the commutator of Hamiltonian vector fields is Hamiltonian. Hence, each $f_{ij}^iX_i+f_{ij}^jX_j$ is a Hamiltonian vector field and it admits a certain Hamiltonian function $\Upsilon_{ij}$, i.e. $\iota_{f_{ij}^iX_i+f_{ij}^jX_j}
\omega=d\Upsilon_{ij}$. Then,
$$
d^2\Upsilon_{ij}=d\iota_{f_{ij}^iX_i+f_{ij}^jX_j}\omega=0,\qquad 1\leq i<j\leq r.
$$
Hence,
\begin{equation}\label{Eq:Con}
0=d(f_{ij}^idh_i+f_{ij}^jdh_j)=df_{ij}^i\wedge dh_i+df_{ij}^j\wedge dh_j,\qquad 1\leq i<j\leq r.
\end{equation}
Since $X_1\wedge\ldots\wedge X_r$ is not vanishing and the mapping $\omega^\flat:v\in TN\mapsto \omega(v,\cdot)\in T^*N$ is an isomorphism because $\omega$ is symplectic and therefore nondegenerate, one has that $dh_1\wedge\ldots\wedge dh_r\neq 0$. Then, $h_1,\ldots,h_r$ are functionally independent functions and some additional functions $y_1,\ldots,y_s$ can be added to them so as to obtain a coordinate system on $T^*N$. Using this coordinate system in (\ref{Eq:Con}), it follows that
$$
\sum_{k=1}^r\frac{\partial f_{ij}^i}{\partial h_k}dh_k\wedge dh_i+\sum_{l=1}^s\frac{\partial f_{ij}^i}{\partial y_l}dy_l\wedge dh_i+\sum_{k=1}^r\frac{\partial f_{ij}^j}{\partial h_k}dh_k\wedge dh_j+\sum_{l=1}^s\frac{\partial f_{ij}^j}{\partial y_l}dy_l\wedge dh_j=0,\quad 1\leq i<j\leq r.
$$
Since $i\neq j$, the linear independence of the basis $dh_1,\ldots,dh_r,dy_1,\ldots,dy_s$ allows one to write
$$
\begin{gathered}
\frac{\partial f_{ij}^i}{\partial h_k}dh_k\wedge dh_i=0,\quad \frac{\partial f_{ij}^i}{\partial y_l}dy_l\wedge dh_i=0,\quad \frac{\partial f_{ij}^j}{\partial y_l}dy_l\wedge dh_j=0,\quad \frac{\partial f_{ij}^j}{\partial h_k}dh_k\wedge dh_j=0,\,\,\,k\notin \{i,j\}\\
\left(\frac{\partial f_{ij}^i}{\partial h_j}-\frac{\partial f_{ij}^j}{\partial h_i}\right)dh_i\wedge dh_j=0.
\end{gathered}
$$
The four equalities in the first line above give that $f_{ij}^i=f_{ij}^i(h_i,h_j)$ and $f_{ij}^j=f_{ij}^j(h_i,h_j)$. Moreover, the second line above yields
$$
\frac{\partial f_{ij}^i}{\partial h_j}=\frac{\partial f_{ij}^j}{\partial h_i},\qquad 1\leq i<j\leq r.
$$
Consequently, there exists a series of functions $h_{ij}=h_{ij}(h_i,h_j)$, with $1\leq i<j\leq r$, such that
$$
f_{ij}^i=\frac{\partial h_{ij}}{\partial h_i},\qquad f_{ij}^j=\frac{\partial h_{ij}}{\partial h_j},\qquad 1\leq i<j\leq r.
$$
Let us prove the converse. If $\{h_i,h_j\}=h_{ij}(h_i,h_j)$ for $1\leq i,j\leq r$, then
$$
d\{h_i,h_j\}=\frac{\partial h_{ij}}{\partial h_i}dh_{i}+\frac{\partial h_{ij}}{\partial h_j}dh_{j},\qquad 1\leq i<j\leq r.
$$
Since $d\{h_i,h_j\},h_i,h_j$ are the Hamiltonian functions for $-[X_i,X_j],X_i,X_j$, respectively, the above expression implies that
$$
-[X_i,X_j]=\frac{\partial h_{ij}}{\partial h_i}X_{i}+\frac{\partial h_{ij}}{\partial h_j}X_{j},\qquad 1\leq i<j\leq r,
$$
and the vector fields $X_1,\ldots,X_r$ form a quasi-rectifiable family.

\end{proof}

\end{document}
